% SEGMENT OLP-0713-S01
% Part: proof-theory
% Chapter: sequent-calculus
% Section: translations

\documentclass[../../../include/open-logic-section]{subfiles}

\begin{document}

\olfileid{pt}{seq}{tr}

\olsection{\Log{G1c} மற்றும் \Log{G3c} இடையிலான மாற்றங்கள்}

\Log{G1c}, \Log{G3c} ஆகிய இரண்டும் செவ்வியல் தருக்கத்துக்கு
வலிதன்மையும் நிறைவுத்தன்மையும் உடையவை என்று முன்பு குறிப்பிட்டோம்.
அதாவது, ஒவ்வொரு !!{structure} இலும் மெய்யான
தொடரணிகளை மட்டுமே, அவை அனைத்தையும், ஒவ்வொரு முறைமையும்
நிறுவுகிறது. குறிப்பாக, இரண்டும் ஒரே தொடரணிகளை நிறுவுகின்றன.
இப்போது வலிதன்மை, நிறைவுத்தன்மை தேற்றங்கள் வழியாகச்
செல்லாமல், நிறுவல் கோட்பாட்டுக்குள்ளேயே இதை நிறுவலாம்.
இந்த வகையான நிறுவல், \Log{G1c} இல் உள்ள
!!{proof}s ஐ \Log{G3c} இல் உள்ள !!{proof}s ஆகவும்
மறுதிசையிலும் மாற்றும் \emph{மாற்றங்களை} வழங்குகிறது.

% SEGMENT OLP-0713-S02
\begin{prop}
$\Log{G1c} \Proves S$ எனில் $\Log{G3c} \Proves S$.
\end{prop}

\begin{proof}
  $\pi$ என்பது \Log{G1c} இல் !!a{proof} எனில்,
  $S$ க்கு \Log{G3c} இல் !!a{proof}~$\pi'$ உண்டு
  என்று காட்டுவோம். இதற்கு $n = \pheight{\pi}$
  மீது தொகுத்தறிதலைப் பயன்படுத்துகிறோம்.

% SEGMENT OLP-0713-S03
  $n = 0$ எனில் $\pi$ ஓர் அடிக்கோள்.
  $\lfalse \Sequent \quad$ என்பது \Log{G3c}
  இலும் ஓர் அடிக்கோள்; இங்கு சூழல்கள் $\Gamma,
  \Delta$ இரண்டும் வெறுமையாக எடுக்கப்படுகின்றன.
  பொதுவான, அணுவாக இருக்க வேண்டியதில்லாத $!A$ உடைய
  $!A \Sequent !A$ என்பது \Log{G3c} இன்
  அடிக்கோளாக வேண்டியதில்லை. ஆனால் இத்தகைய
  எல்லாத் தொடரணிகளும் \Log{G3c} இல்
  !!{provable} என்பதை
  \olref[ex]{prop:G3c-axioms} இல் நிறுவியுள்ளோம்.

% SEGMENT OLP-0713-S04
  $n>0$ எனில் $\pi$ ஒரு விதி~$R$ இல் முடிகிறது.
  தொகுத்தறி கருதுகோளின்படி அந்த விதியின்
  முன்கூற்றுகளுக்கு \Log{G3c} இல் !!{proof}s
  உள்ளன; அதாவது உடனடி உள்-!!{proof}s க்கு
  \Log{G3c} இல் மாற்றங்கள் உள்ளன.
  $R$ இரு முறைமைகளிலும் ஒரே வடிவிலிருந்தால்,
  மொழிமாற்றப்பட்ட முன்கூற்றுகளின் !!{proof}s
  மீது அந்த விதியை நேரடியாகப் பயன்படுத்தலாம்.
  இட இணைப்பு, வல பிரிப்பு, இட அனைத்தளவையடை,
  வல இருப்பளவையடை விதிகளின் \Log{G3c}
  முன்கூற்றில் கூடுதல் வாய்பாடு தேவைப்படும்;
  \olref[adm]{prop:weak-G3c-adm} இன் வலுவிழக்கச்
  சேர்த்தல் அனுமதியால் அதைச் சேர்த்தபின் அந்த
  முறைமையின் தொடர்புடைய விதியைப் பயன்படுத்துகிறோம்.
  குறிப்பாக, $R$ என்பது $\LeftR{\land}$ அல்லது
  $\RightR{\lor}$ எனில் இந்த முறையே பொருந்தும்.
  வலுவிழக்கச் சேர்த்தல் விதி எனில்
  \olref[adm]{prop:weak-G3c-adm} ஐப் பயன்படுத்துகிறோம்;
  சுருக்க விதி எனில் \olref[inv]{prop:G3c-cont-adm}
  ஐப் பயன்படுத்துகிறோம்.
\end{proof}

% SEGMENT OLP-0713-S05
\begin{prop}
$\Log{G3c} \Proves S$ எனில் $\Log{G1c} \Proves S$.
\end{prop}

\begin{proof}
  மீண்டும், $S$ க்கான !!a{proof}~$\pi$ இன்
  !!{height} மீது தொகுத்தறிதலால் செயல்படுவோம்.
  \Log{G3c} இன் ஒவ்வோர் அடிக்கோளையும்
  \Log{G1c} இல் ஓர் அடிக்கோளிலிருந்து
  \LeftR{\Weakening}, \RightR{\Weakening}
  விதிகளைப் பலமுறை பயன்படுத்தி !!{prove} செய்யலாம்:
  \[
  \Axiom$!C \fCenter !C$
  \RightLabel{\LeftR{\Weakening}}
  \Deduce$!C, \Gamma \fCenter !C$
  \RightLabel{\RightR{\Weakening}}
  \Deduce$!C, \Gamma \fCenter \Delta, !C$
  \DisplayProof
\qquad
  \Axiom$\lfalse \fCenter$
  \RightLabel{\LeftR{\Weakening}}
  \Deduce$\lfalse, \Gamma \fCenter$
  \RightLabel{\RightR{\Weakening}}
  \Deduce$\lfalse, \Gamma \fCenter \Delta$
  \DisplayProof
  \]

% SEGMENT OLP-0713-S06
  $\pi$ ஒரு விதி~$R$ இல் முடிந்தால்,
  தொகுத்தறி கருதுகோள் அதன் முன்கூற்றுகளுக்கு
  \Log{G1c} இல் !!{proof}s தருகிறது.
  இரு முறைமைகளிலும் வடிவம் ஒன்றாக உள்ள
  விதி~$R$ ஐ அந்த !!{proof}s மீது
  அப்படியே பயன்படுத்தலாம். ஆனால்
  \LeftR{\land}, \RightR{\lor} விதிகள்
  \Log{G1c}, \Log{G3c} இல் வேறுபடுகின்றன.
  \Log{G3c} இன் வடிவங்களை \Log{G1c} இல்
  பின்வருமாறு வருவிக்கலாம்:
  \[
  \Axiom$!A, !B, \Gamma \fCenter \Delta$
  \RightLabel{\LeftR{\land}}
  \UnaryInf$!A \land !B, !B, \Gamma \fCenter \Delta$
  \RightLabel{\LeftR{\land}}
  \UnaryInf$!A \land !B, !A \land !B, \Gamma \fCenter \Delta$
  \RightLabel{\LeftR{\Contraction}}
  \UnaryInf$!A \land !B, \Gamma \fCenter \Delta$
    \DisplayProof
\qquad
  \Axiom$\Gamma \fCenter \Delta, !A, !B$
  \RightLabel{\RightR{\lor}}
  \UnaryInf$\Gamma \fCenter \Delta, !A \lor !B, !B$
  \RightLabel{\RightR{\lor}}
  \UnaryInf$\Gamma \fCenter \Delta, !A \lor !B, !A \lor !B$
  \RightLabel{\RightR{\Contraction}}
  \UnaryInf$\Gamma \fCenter \Delta, !A \lor !B$
    \DisplayProof
  \]
  இட அனைத்தளவையடை, வல இருப்பளவையடை
  விதிகளின் முன்கூற்றுகளில்
  முதன்மை வாய்பாடும் தக்கவைக்கப்படுகிறது.
  \Log{G1c} இன் தொடர்புடைய விதியை
  அந்தத் தக்கவைக்கப்பட்ட வாய்பாட்டைச் சூழலாகக்
  கொண்டு பயன்படுத்தி, முடிவில் கிடைக்கும்
  இரட்டைப் பிரதிகளை இட அல்லது வல சுருக்கத்தால்
  ஒன்றாக்குகிறோம்.
\end{proof}

\end{document}
